\begin{figure}
\centering
\small
\begin{tabular}{>{\raggedright\arraybackslash}p{0.23\textwidth}>{\raggedright\arraybackslash}p{0.55\textwidth}>{\raggedright\arraybackslash}p{0.12\textwidth}}
\hline
\textbf{Python error} & \textbf{Ruled out when} & \textbf{See} \\
\hline
\pyinline{NameError}, \pyinline{UnboundLocalError} & Always: a variable may be referenced only where it is definitely assigned. & \figref{well-formed}, \figref{well-formed-expr} \\[1mm]
\pyinline{ModuleNotFoundError} & Always: an import requires its target to be a module of the program. & \figref{well-formed-imports} \\[1mm]
\pyinline{ImportError} & Always: a from-import requires each imported name to be a member of the source module, and import cycles are rejected. (A Python cycle may run or fail, depending on where the cycle is entered.) & \figref{imports}, \figref{well-formed-module} \\[1mm]
\pyinline{AttributeError} & The reference is a module attribute, which must resolve to a definitely assigned member.  & \figref{well-formed-expr} \\[1mm]
\pyinline{TypeError} & The operation is a constructor call or class pattern, which must match the class's arity exactly. (Only excess sub-patterns are a Python error; a pattern with fewer is excluded even though Python permits it.) & \figref{well-formed-expr}, \figref{well-formed-pat} \\
\hline
\end{tabular}
\caption{Python errors prevented by well-formedness}
\label{fig:prevented-errors}
\end{figure}
